\newpage
\section{Déroulement du projet}
\label{sec:deroulement}

\subsection{Réalisation du projet par étapes}
\paragraph{} Le projet a été réalisé selon les étapes suivantes :

\begin{enumerate}
\item \textbf{Semaine 1-2} : Mise en place des classes de base (Card, Deck, Player)
\item \textbf{Semaine 3-4} : Implémentation des combinaisons (Combination, determineType, canBeat)
\item \textbf{Semaine 5-6} : Développement des bots et de la logique de jeu (Game)
\item \textbf{Semaine 7-8} : Interface graphique (Swing) et refonte du layout
\item \textbf{Semaine 9-10} : Tests unitaires (JUnit) et corrections de bugs
\item \textbf{Semaine 11-12} : Rédaction du rapport et préparation de la livraison
\end{enumerate}

\subsection{Répartition des tâches, communication et intégration}
\paragraph{} La répartition des tâches s'est organisée comme suit :

\begin{itemize}
\item \textbf{Laiza KHODJA} : Interface graphique (GameStyle, MusicPlayer, EndGameGUI, BotsPanel)
\item \textbf{Nesrine BOUATMANE} : Classes de données (Card, Deck, Player, GameStats, tests unitaires)
\item \textbf{Nassim BERRANDOU} : Logique métier (Game, Combination, bots, correction bugs)
\end{itemize}

\paragraph{} La communication s'est faite via Discord et GitHub. Le versionnement Git a permis une intégration continue des différentes parties.

\paragraph{} Les principaux défis rencontrés ont été :
\begin{itemize}
\item La gestion du reset du tas et du recyclage des cartes
\item La compatibilité de la musique entre Linux et Windows
\item La synchronisation des tours entre le joueur humain et les bots
\end{itemize}
